\begin{theorem}[sigmaR-eq-sum-div]
\uses{sigmaR}
\label{sigmaR_eq_sum_div}
\lean{ArithmeticFunction.sigmaR_eq_sum_div}
% at /home/runner/work/trough/trough/library/PrimeNumberTheoremAnd/IwaniecKowalskiCh1.lean:439.0-455.99
\leanok
We have that $\sigma^R_s(n)=\sum_{d\mid n}(n/d)^s$.
\end{theorem}
\begin{proof}
\uses{sigmaR_apply}
\leanok
Note that $d \mapsto n/d$ forms a one-to-one mapping between the divisors of $n$. Using this in
    combination with the definiton we have that
    $$\sigma^R_s(n)=\sum_{d\mid n}d^s=\sum_{d\mid n}(n/d)^s.$$
\end{proof}
